\documentclass[11pt]{amsart}
\usepackage[margin=1.2in]{geometry}
\usepackage{amsmath,amssymb,amsthm}
\usepackage{xcolor}
\usepackage[colorlinks=true,linkcolor=blue!60!black,citecolor=blue!60!black,urlcolor=blue!60!black]{hyperref}

\newtheorem{theorem}{Theorem}
\newtheorem{proposition}[theorem]{Proposition}
\newtheorem{lemma}[theorem]{Lemma}
\newtheorem{corollary}[theorem]{Corollary}
\theoremstyle{remark}
\newtheorem{remark}[theorem]{Remark}

\newcommand{\RR}{\mathbb{R}}
\newcommand{\FF}{\mathbb{F}}
\newcommand{\Gb}{\overline{\mathcal{G}}}
\newcommand{\Gm}{\mathcal{G}}
\newcommand{\Gu}{\underline{\mathcal{G}}}
\newcommand{\Gbar}{\overline{G}}
\newcommand{\Stab}{\operatorname{Stab}}

\title[Mutation graphs of uniform oriented matroids on $\le 9$ elements]{The mutation graphs of uniform oriented matroids\\ on at most nine elements are connected}
\author{Reuel Lee}
\address{}
\email{reuellee@gmail.com}
\date{July 31, 2026}

\begin{document}

\begin{abstract}
Cordovil and Las Vergnas conjectured that the mutation graph of uniform
oriented matroids of fixed rank and size is connected. Knauer and Marc
distinguish three versions of this graph --- on labelled uniform oriented
matroids, on their reorientation classes, and on their isomorphism
classes --- prove the labelled graph connected in rank $3$ for every $n$,
prove the conjecture for $n\le 9$ at the reorientation-class level, and
write that beyond rank $3$ they know nothing about the labelled graph,
where they suspect a counterexample exists. We show there is no
counterexample below $n=10$: for every rank $r$ and every $n\le 9$, all
three graphs are connected. Once their rank-$3$ theorem, oriented-matroid
duality and the elementary ranks ($r\le2$ and $r\ge n-1$) are accounted
for, the cases genuinely settled here are $(n,r)=(8,4)$, $(9,4)$ and its
dual $(9,5)$ --- among them the case they single out as hardest. The
method is a covering-space reduction expressing the number of components
of the labelled graph as the index of a holonomy subgroup computable on
the isomorphism-class quotient. Two by-products: the number of
isomorphism classes of uniform rank-$4$ oriented matroids on $9$ elements
is $9{,}276{,}595$, proved by a mass identity against a structurally
independent extension count, correcting the value $9{,}276{,}601$ printed
by Fukuda, Miyata and Moriyama; and, when the isomorphism-class graph is
connected and the holonomy surjects onto $S_n$, for odd $n$ the number of
components of the labelled graph is forced into $\{1,2^{n-1}\}$ with no
intermediate value possible. All computations are exact. Certificates and
standalone checkers, all of them in the repository, accompany the holonomy
claims, the $9{,}276{,}595$-class enumeration, and a mutation spanning tree
joining those classes into a single component; \S\ref{sec:honesty} states
precisely which claims they cover and which remain reproducible rather
than certified.
\end{abstract}

\maketitle

\section{Introduction}

Let $\mathcal{M}$ be a uniform oriented matroid (UOM) of rank $r$ on
$E=\{1,\dots,n\}$, presented by a chirotope $\chi$. A \emph{mutation} of
$\mathcal{M}$ flips the sign of a single basis so that the result is
again a chirotope; geometrically, for realizable $\mathcal{M}$, it is the
passage of one point through the hyperplane spanned by $r-1$ others. The
\emph{mutation graph} has the UOMs of rank $r$ on $n$ elements as
vertices and mutations as edges.

Cordovil and Las Vergnas conjectured that this graph is connected; the
conjecture is recorded in \cite{RS88}. Knauer and Marc \cite{KM23} observe that the
statement depends on which objects one takes as vertices, and separate
three graphs, in decreasing fineness:
\[
  \Gb^{n,r} \quad(\text{labelled UOMs}), \qquad
  \Gm^{n,r} \quad(\text{reorientation classes}), \qquad
  \Gu^{n,r} \quad(\text{isomorphism classes}),
\]
where reorientation means the action of $\{\pm1\}^n$ (sign changes) and
isomorphism the action of the full hyperoctahedral group
$\{\pm1\}^n\rtimes S_n$ (sign changes and relabelling). Connectivity of
each implies connectivity of the next. Knauer and Marc state the
Cordovil--Las Vergnas conjecture at the middle level. They prove, at the
\emph{labelled} level, that $\Gb^{n,3}$ is connected for every $n$
\cite[Prop.~3.2]{KM23}, via Ringel's homotopy theorem together with an
explicit reorientation gadget; and they deduce $\Gm^{n,r}$ connected for
$n\le 9$ from a computation that $\Gu^{n,r}$ is connected in that range.
Of the finest level in higher rank they write that checking connectivity
``is far more demanding'' and that they ``do not know anything about the
connectivity of this graph beyond rank~3''; in their conclusion they say
they ``suspect the existence of a counter example at least in the setting
of $\Gb^{n,r}$.''

Our main result is that no such counterexample exists below $n=10$.

\begin{theorem}\label{thm:main}
For every $n\le 9$ and every rank $r$, the graphs $\Gb^{n,r}$,
$\Gm^{n,r}$ and $\Gu^{n,r}$ are connected.
\end{theorem}

The content is the labelled level; the other two follow. Much of the
range is not new, and it is worth saying at once which part is. Ranks
$r\le 1$ and $r\ge n-1$ are degenerate and elementary
(\S\ref{sec:mass}); rank $2$ is elementary as well; rank $3$ is
\cite[Prop.~3.2]{KM23}, and rank $n-3$ follows from it by
oriented-matroid duality. What remains at $n\le 9$ is rank $4$ for
$n=8,9$ together with $(9,5)$, which is dual to $(9,4)$. So the cases
genuinely settled here are
\[
  (n,r)\;=\;(8,4),\qquad (9,4),\qquad (9,5)\cong(9,4)^{*},
\]
and it is $(9,4)$ that carries the weight: its isomorphism-class
quotient has $9{,}276{,}595$ vertices — the case Knauer and Marc, who
computed its quotient level, single out as the computationally most
demanding — and the labelled level there is what they describe as
beyond present knowledge. Theorem~\ref{thm:main} does not settle the
Cordovil--Las Vergnas conjecture, nor the labelled question, in general:
it rules out a counterexample only for $n\le 9$.

Two by-products deserve separate statement. The first corrects the
literature.

\begin{theorem}\label{thm:count}
The number of isomorphism classes of uniform oriented matroids of rank
$4$ on $9$ elements is exactly $9{,}276{,}595$.
\end{theorem}

This value appears in Finschi's database \cite{Finschi} and in
\cite[Table 1]{KM23}; Fukuda, Miyata and Moriyama \cite[Table]{FMM13}
print $9{,}276{,}601$, which is six too many. Our argument is not a third
opinion but an exact integer identity: the orbit masses of the classes
found by an exhaustive search sum to a target obtained by a structurally
independent route, a single-element extension sweep over the $2{,}628$
classes at $n=8$. Under hypotheses made explicit in \S\ref{sec:mass} ---
the enumerated classes valid, pairwise inequivalent, and carrying their
true stabilizer orders, and the target exact --- a search that missed a
class falls short and cannot overshoot. The first three hypotheses are
discharged by a standalone checker over all $9{,}276{,}595$ classes; the
fourth is reproducible but not independently certified
(\S\ref{sec:honesty}).

The second by-product explains why the labelled question at odd $n$ is
binary.

\begin{proposition}\label{prop:dichotomy}
Let $n$ be odd, let $\Gu^{n,r}$ be connected, and suppose the holonomy
subgroup $H$ of \S\ref{sec:reduction} surjects onto $S_n$. Then
\[
  \#\mathrm{components}\bigl(\Gb^{n,r}\bigr)\in\{1,\ 2^{n-1}\}.
\]
For $n=9$: either $1$ or $256$, with no intermediate value possible.
\end{proposition}

Finally, we record an erratum: \cite[Table 1]{KM23} lists $482$ classes at
$(n,r)=(9,3)$. The correct value is $4382$, as its own dual entry at
$(9,6)$, Finschi's database, \cite{FMM13} and our independent generation
all agree.

\subsection*{Verifying this note before reading it}
The computational claims are machine-checkable, within the scope
\S\ref{sec:honesty} sets out. From the repository
\begin{center}
\url{https://github.com/reuellee/finite-certificates}
\end{center}
(directory \texttt{ai/omgamma}), with Python 3 and no third-party
packages beyond \texttt{numpy} for the fast checker:
\begin{verbatim}
    python verify_omgamma.py             # holonomy certificates
    python verify_omgamma.py --canary    # self-test: must fail on a
                                         # missing artifact
    python submodules.py                 # Proposition 3, two methods
    python coverage_checker.py --artifact data/coverage_4_9 --workers 2 \
           --extcount data/extcount_4_9.jsonl  # coverage + reachability,
                                               # all 9.28M classes
    python coverage_checker.py --canary --artifact data/coverage_4_9
                                               # 12 sabotages; must REJECT
\end{verbatim}
Every artifact these commands read is in the repository, the coverage
certificate included: it is a $10.4$~MB file from which the checker
reconstructs all $9{,}276{,}595$ class keys and stabilizer orders.
\S\ref{sec:honesty} says how far each command reaches.
The ledger \texttt{ai/omgamma/OMGAMMA.md} records every computation, its
artifacts, its trust boundaries, and its errata --- including one of our
own, discussed in \S\ref{sec:honesty}.

\section{The three graphs}\label{sec:defs}

Fix $n$ and $r$. A uniform chirotope is an alternating map
$\chi\colon E^r\to\{\pm1\}$ satisfying the three-term Grassmann--Pl\"ucker
sign conditions; $\chi$ and $-\chi$ describe the same oriented matroid,
so a vertex of $\Gb^{n,r}$ is the pair $\{\chi,-\chi\}$. Let
\[
  G' \;=\; S_n\times\{0,1\}^n\times\{0,1\},
\]
\[
  \bigl((\sigma,\varepsilon,s)\cdot\chi\bigr)(x_1,\dots,x_r)
   \;=\;(-1)^{s}\,(-1)^{\sum_{i=1}^{r}\varepsilon_{x_i}}\,
        \chi(\sigma^{-1}x_1,\dots,\sigma^{-1}x_r),
\]
where $\varepsilon=(\varepsilon_1,\dots,\varepsilon_n)\in\{0,1\}^n$
records which ground-set elements are reoriented and $s$ is the global
sign. Here $G'$ carries the semidirect-product multiplication induced by
the $S_n$-action on $\{0,1\}^n$,
\[
  (\sigma_1,\varepsilon_1,s_1)(\sigma_2,\varepsilon_2,s_2)
  =(\sigma_1\sigma_2,\ \varepsilon_1\oplus\sigma_1(\varepsilon_2),\ s_1\oplus s_2).
\]
\sloppy
Throughout Sections~\ref{sec:reduction}--\ref{sec:mass} we work in the
nondegenerate range $2\le r\le n-2$; the degenerate ranks are disposed of
directly in Section~\ref{sec:mass}. There the subgroup
$K_4=\{(\mathrm{id},0,0),(\mathrm{id},1^n,0),(\mathrm{id},0,1),
(\mathrm{id},1^n,1)\}$ acts trivially on pairs $\{\chi,-\chi\}$, and we
write $\Gbar=G'/K_4$, of order $n!\,2^{\,n-1}$, with normal sign subgroup
$\bar R=\{0,1\}^n/\langle 1^n\rangle$ and quotient
$\pi\colon\Gbar\twoheadrightarrow S_n$. (Nothing below requires $K_4$ to
be the \emph{full} kernel; it is, in this range, but we use only that it
acts trivially.) Then
\[
  \Gm^{n,r}=\Gb^{n,r}/\bar R, \qquad \Gu^{n,r}=\Gb^{n,r}/\Gbar ,
\]
and mutation is $\Gbar$-equivariant, so these quotients are graphs in the
usual way.

\section{From the labelled graph to a subgroup}\label{sec:reduction}

The reduction is standard covering-space bookkeeping, stated for a group
action that is \emph{not} free --- which is exactly where care is needed.

\begin{lemma}\label{lem:transitive}
Let a finite group $\Gbar$ act by automorphisms on a graph $\Gb$ with
connected quotient $\Gu=\Gb/\Gbar$. Then $\Gbar$ permutes the components
of $\Gb$ transitively, every component maps onto $\Gu$, and for a base
vertex $v_0$,
\[
  \#\mathrm{components}(\Gb)=[\Gbar:H],\qquad
  H:=\{g\in\Gbar\;:\;g\cdot v_0\in\mathrm{comp}(v_0)\}.
\]
\end{lemma}

For a chirotope $\chi$ and a basis $B$ at which $\chi$ is mutable, write
$\mu_B\chi$ for the chirotope agreeing with $\chi$ except that its sign
on $B$ is negated. If the class of $\mu_B\chi_c$ is $c'$, the
\emph{voltage} of that mutation edge is the element $t\in\Gbar$ with
$\mu_B\chi_c=t\cdot\chi_{c'}$ (as pairs); it is well defined up to
$\Stab_{\Gbar}(\chi_{c'})$, and any choice may be fixed.

\begin{lemma}\label{lem:gens}
Fix class representatives $\chi_c$ and a spanning tree of $\Gu$ rooted at
$c_0$. Define transports $\tau_c\in\Gbar$ by $\tau_{c_0}=\mathrm{id}$
and, along each tree edge directed away from the root, $c\to c'$ with
voltage $t$, by $\tau_{c'}=\tau_c\,t$. Then $H$ is generated by
\[
  \underbrace{\{\tau_c\,u\,\tau_c^{-1}\;:\;u\in\Stab_{\Gbar}(\chi_c)\}}_{\text{(i) stabilizers}}
  \;\cup\;
  \underbrace{\{\tau_c\,t\,\tau_{c'}^{-1}\;:\;\text{every mutation edge }(c,j,c')\}}_{\text{(ii) voltages}} .
\]
\end{lemma}

\begin{proof}
$(\supseteq)$ Induction along the tree gives $\tau_c\cdot\chi_c\in
\mathrm{comp}(v_0)$ for every $c$, mutations being $\Gbar$-equivariant.
Then $\tau_c u\tau_c^{-1}$ fixes the component vertex $\tau_c\cdot\chi_c$,
and $\tau_c t\tau_{c'}^{-1}$ carries $\tau_{c'}\cdot\chi_{c'}$ to
$\tau_c\cdot(\mu_B\chi_c)$, a neighbour of $\tau_c\cdot\chi_c$; an
element carrying some component vertex into the component stabilizes it.
Note this direction needs only that the tree paths exist, not that the
search is complete, so elements harvested from a partial search are
genuine members of $H$.

$(\subseteq)$ Let $g\in H$, witnessed by a walk
$v_0=x_0\sim x_1\sim\dots\sim x_k=g\cdot v_0$ in $\Gb$. Write
$x_i=s_i\cdot\chi_{c_i}$ with $s_0=\mathrm{id}$. Applying $s_i^{-1}$ to
the edge $x_i\sim x_{i+1}$ shows $s_i^{-1}s_{i+1}\cdot\chi_{c_{i+1}}$ is
a mutation of $\chi_{c_i}$, so
$s_i^{-1}s_{i+1}\in t_i\Stab_{\Gbar}(\chi_{c_{i+1}})$ for the voltage
$t_i$ of some mutation edge $c_i\to c_{i+1}$. Hence $g=s_k h$ with
$h\in\Stab_{\Gbar}(\chi_{c_0})$ and $s_k$ a product of the $s_i^{-1}s_{i+1}$;
inserting $\tau_{c_i}^{-1}\tau_{c_i}$ between consecutive factors
rewrites that product as a word in the displayed conjugated families,
which telescopes because $\tau_{c_0}=\mathrm{id}$.
\end{proof}

Family (i) is not optional: the action of $\Gbar$ on chirotopes has
nontrivial stabilizers, and omitting the stabilizer conjugates computes a
proper subgroup of $H$. We learned this the hard way
(\S\ref{sec:honesty}).

Combining, with $\Gu$ connected: $\Gb^{n,r}$ is connected iff
$H=\Gbar$. For the middle level, the components of
$\Gm^{n,r}=\Gb^{n,r}/\bar R$ are the $\bar R$-orbits on the components of
$\Gb^{n,r}$; since $\bar R$ is normal, $\bar R H$ is a subgroup, and
Lemma~\ref{lem:transitive} gives
\[
  \#\mathrm{components}(\Gm^{n,r})=[\Gbar:\bar R H]=[S_n:\pi(H)] .
\]
So the labelled question becomes: compute a subgroup of $\Gbar$ from data
harvested on the quotient, and ask whether it is everything.

\section{The dichotomy at odd $n$}\label{sec:dichotomy}

Write $S:=H\cap\bar R$ for the sign part of $H$. Since $\bar R$ is normal
and abelian, conjugation gives
$(\sigma,\delta,s)(\mathrm{id},\varepsilon,0)(\sigma,\delta,s)^{-1}
 =(\mathrm{id},\sigma(\varepsilon),0)$, so $S$ is $\pi(H)$-invariant. If
$\pi(H)=S_n$ then $S$ is an $\FF_2[S_n]$-submodule of $\bar R$.

\begin{proof}[Proof of Proposition~\ref{prop:dichotomy}]
For odd $n$ we have $1^n\notin E$, where $E\le\FF_2^n$ is the even-weight
subspace, so $\FF_2^n=\langle 1^n\rangle\oplus E$ and $\bar R\cong E$. We
check directly that $E$ is an irreducible $\FF_2[S_n]$-module. Let
$0\ne V\le E$ be invariant and $0\ne v\in V$. Since $v$ has even weight
and $n$ is odd, $v$ has a coordinate $i$ in its support and a coordinate
$j$ outside it; then $v+(i\,j)\cdot v=e_i+e_j$ lies in $V$. The
$S_n$-orbit of $e_i+e_j$ is all weight-two vectors, and these span $E$,
so $V=E$. Hence $S\in\{0,\bar R\}$, and by Lemma~\ref{lem:transitive} the
number of components is $[\Gbar:H]\in\{2^{n-1},1\}$.
\end{proof}

The conclusion is corroborated computationally, independently of the
argument: \texttt{submodules.py} enumerates the $S_n$-invariant subspaces
of $\FF_2^n$ for $n=5,\dots,10$ as sums of spans of weight classes,
finding exactly four for each $n$; for $n\le 7$ it also enumerates
\emph{all} subspaces through their reduced row echelon forms and filters
them, the totals $374/2825/29212$ matching the Galois numbers (so that
enumeration is provably exhaustive) and agreeing with the first method
where both were run. A deliberately non-invariant subspace is rejected by
the invariance test. For even $n$ an intermediate value is admissible;
the full sign parts observed at $n=8$ are therefore not forced by
algebra.

Proposition~\ref{prop:dichotomy} is not needed for the final verdict ---
the certificate exhibits a full sign span directly --- but it explains
the shape of the computation: at $n=9$ the tracked sign dimension jumped
from its minimum to its maximum with no intermediate value, and the
proposition says no intermediate value was ever available.

\section{The computation}\label{sec:mass}

Ranks $r\le 1$ and $r\ge n-1$ are degenerate: the three-term
Grassmann--Pl\"ucker condition set is empty (it is nonempty only when
$r\ge2$ and $n\ge r+2$), so every sign vector on the $\binom{n}{r}$ bases
is a chirotope and every single-bit flip is a mutation; $\Gb^{n,r}$ is
then a folded cube, visibly connected. Rank $2$ is likewise elementary: a
uniform oriented matroid of rank $2$ on $n$ elements is a cyclic
arrangement of $n$ signed points on a circle, so its chirotope is
determined by that cyclic order together with the reorientation, a
mutation transposes two neighbours, and adjacent transpositions connect
all such data. Rank $3$ is \cite[Prop.~3.2]{KM23}. Rank $4$ is computed
exhaustively as described below. Ranks $5\le r\le n-2$ follow by duality. For a basis
$B=\{x_1<\dots<x_r\}$ with complement
$E\setminus B=\{x_{r+1}<\dots<x_n\}$, the dual chirotope is
\[
  \chi^*(x_{r+1},\dots,x_n)\;=\;\chi(x_1,\dots,x_r)\cdot
    \operatorname{sgn}(x_1,\dots,x_n),
\]
$\operatorname{sgn}$ being the sign of the permutation
$(x_1,\dots,x_n)$ of $E$. Then $\chi\mapsto\chi^*$ is a graph
isomorphism $\Gb^{n,r}\cong\Gb^{n,n-r}$ (and likewise for $\Gm$ and
$\Gu$), carrying the mutation at $B$ to the mutation at $E\setminus B$:
since $\chi^*(E\setminus B)=\pm\chi(B)$ with a sign depending only on
$B$, negating $\chi$ at $B$ is the same operation as negating $\chi^*$ at
$E\setminus B$, and validity is preserved on both sides. (This was also
checked computationally, on all $135$ classes at $(8,3)$ among others.)
This covers every rank at $n\le 9$.

For rank $4$ at $n\le 8$ (and, as a control, rank $3$ at $n\le 9$, where
the answer is already known from \cite[Prop.~3.2]{KM23}) the quotient
$\Gu$ is small enough for an exhaustive breadth-first search harvesting
both generator families of Lemma~\ref{lem:gens} at every class. In each
case $H=\Gbar$, so all three graphs are connected; the rank-$3$ runs
agreeing with \cite{KM23} is one of this project's regression controls.

The case $(9,4)$ needs $9{,}276{,}595$ classes and
$150{,}561{,}898$ directed edge traversals, run as a disk-based
level-synchronous search. Completeness is certified by a \emph{mass
identity}. Every class $c$ contributes an orbit of size
$|G'_9|/|\Stab(\chi_c)|$ labelled objects, and the total number of
labelled uniform chirotopes is computed by a structurally independent
counting route --- counting single-element extensions of each of the
$2{,}628$ classes at $(8,4)$, which touches neither the mutation graph
nor the search:
\[
  \sum_{c\ \text{found}}\frac{|G'_9|}{|\Stab(\chi_c)|}
  \;=\;1{,}722{,}704{,}635{,}330{,}560
  \;=\;N_{\chi}(4,9),
\]
an exact integer identity, reached exactly. Provided the enumerated
classes are valid, pairwise inequivalent, and carry their true stabilizer
orders, and provided the extension target is exact, a missed class leaves
the sum short and it cannot overshoot; under those hypotheses the
identity proves that the search visited every class, which is
Theorem~\ref{thm:count}. Those hypotheses are conditions on this
project's own canonicalizer and extension counter, not on the search: a
defective list carrying duplicates or wrong stabilizer orders could in
principle overshoot. The first three are discharged by the coverage
certificate of \S\ref{sec:honesty}, checked by a program sharing no code
with the generator; the fourth is not. Of the classes found, $8{,}913$
have a nontrivial stabilizer.

Completeness of the list is not by itself connectivity of $\Gu^{9,4}$:
those are different statements, and the second also needs the listed
classes to be joined to one another. That is a separate artifact, a
mutation \emph{spanning tree} over all $9{,}276{,}595$ classes
(\S\ref{sec:honesty}), recording for each class but a fixed root the
parent class and the basis mutated to reach it. The same
standalone checker verifies in exact arithmetic that every recorded edge
is a genuine mutation edge of $\Gu^{9,4}$ and that the parent pointers
form a tree with one root, so every listed class is joined to that root
by a path of mutations. All $9{,}276{,}595$ of them therefore lie in one
component; the list being complete, $\Gu^{9,4}$ is connected.

The harvested holonomy is the full group: $\pi(H)=S_9$ of order
$362{,}880$ and the sign part full. Notably, on the shipped certificate
the $73$ stabilizer-derived generators are all even permutations and
generate exactly $A_9$, and a single edge voltage of odd permutation part
supplies the remaining coset --- so on that certificate family (i) alone
does not reach $S_9$. (We do not claim this of the complete family (i)
over all classes, which was not enumerated: harvesting stops once $H$ is
full.) With
$\Gu^{9,4}$ connected, Lemma~\ref{lem:transitive} gives
$[\Gbar:H]=1$, so $\Gb^{9,4}$, $\Gm^{9,4}$ and $\Gu^{9,4}$ are connected,
and $(9,5)$ follows by duality. This completes
Theorem~\ref{thm:main}.

\begin{remark}
The computation was interrupted by an out-of-memory kill and resumed from
a checkpoint. A resumed run harvests edge voltages only from levels
expanded after the resume, so the subgroup it computes is a \emph{lower}
bound $H'\le H$. Since $H'=\Gbar$ was reached, $H=\Gbar$ outright: a
lower bound equal to the whole group is the whole group. The direction
matters --- for a negative verdict a lower bound proves nothing, and a
completing pass over the earlier levels would be mandatory. The shipped
summary records the lower-bound status explicitly.
\end{remark}

\section{Certificates, and what is not certified}\label{sec:honesty}

The claims above rest on artifacts rather than on trust in the search
programs, but the two halves rest on different artifacts and to different
depths, and the first boundary below says how.
Each completed case emits representatives, a spanning tree with voltages,
generators, and exhibits; two checkers sharing no code with the generator
--- one pure Python, one \texttt{numpy} --- re-verify the chirotope
conditions, the group action, the mutation identity of every tree edge,
the generator words, and the resulting group orders from scratch. Five
sabotage mutations (corrupted representative, corrupted voltage,
corrupted generator, truncated exhibits, re-pointed tree parent) are
rejected, and the verifier fails loudly when a required artifact is
missing, renamed or emptied.

What follows states the boundaries plainly, beginning with the one that is
sharp enough to belong before the claims it qualifies.

\emph{The compact certificates prove $H=\Gbar$; the catalog and its
reachability are certified by a separate artifact, and only on one side of
the mass identity.} The compact $(9,4)$ certificates contain the root-path
closure of the classes referenced by generators --- a few hundred classes
--- which is what the holonomy claim needs and all that the checkers above
verify. The left-hand side of the mass identity is a different object: the
$9{,}276{,}595$ classes themselves, with their stabilizer orders and with
mutation edges joining them. That object is now a $10.4$~MB file in the
repository, and it records only
\begin{quote}
the key of one root class --- a $126$-bit integer whose bits are the signs
of a chirotope on the bases in colex order --- and, for each other class,
the class it was reached from together with the colex index $j_c$ of the
basis mutated to reach it.
\end{quote}
There are no class keys in it, no stabilizer orders, no voltages and no
depths: those are \emph{derived}. What makes that possible is an identity
that had been read as a check,
\[
  (\sigma_c,\varepsilon_c,s_c)\cdot\chi_c
   \;=\;\mu_{B_{j_c}}\bigl(\chi_{\mathrm{parent}(c)}\bigr)
   \qquad\text{for some }(\sigma_c,\varepsilon_c,s_c)\in G' .
\]
It says that $\chi_c$ lies in the $G'$-orbit of
$\mu_{B_{j_c}}(\chi_{\mathrm{parent}(c)})$; and the class key is by
construction a function of the orbit alone. Read as a derivation rather
than as a check, it therefore determines each class's key from its
parent's. The rows are stored in a canonical order --- by depth, then by
parent, then by mutated basis --- which is a function of the tree alone,
makes every parent pointer precede its row and the parent array monotone,
hence codable as a $2.3$~MB bitmap, and lets the checker verify the layout
itself. The earlier form of the artifact --- a $62$~MB array of sorted
keys and stabilizer orders beside an $83$~MB witness of parents, mutated
bases, voltages and depths --- is still what the export writes, but it is
redundant, and it is not tracked.

A standalone checker \texttt{coverage\_checker.py}, importing nothing from
this project --- it refuses to run if any module loaded into the
interpreter resolves to a file beside it --- and rebuilding the basis
order, the Grassmann--Pl\"ucker conditions, the sign lattice and the group
action from scratch, reconstructs that catalog and checks it in exact
integer arithmetic. Taking the classes in order, each parent before its
children, it flips the recorded basis of the parent's chirotope and
verifies that the result is a \emph{valid} uniform chirotope --- which is
exactly the statement that $B_{j_c}$ is mutable in the parent, so that the
edge is a real mutation edge --- and then canonicalizes it into the
child's key and stabilizer order, by enumerating every admissible
relabelling and maximising over the sign lattice. It verifies that the
root key it was handed is itself extremal in its orbit; that the
$9{,}276{,}595$ reconstructed keys are pairwise distinct; that every
reconstructed stabilizer order divides $|G'|$ and that their histogram is
the one the manifest declares; and that the orbit masses sum to
$1{,}722{,}704{,}635{,}330{,}560$ with a count of $9{,}276{,}595$. On the
tree it verifies that the packed streams are well formed --- exact bit
lengths, zero padding, one gap bit per non-root class, every basis index
in range
--- and that each parent strictly precedes its row, so that the parent map
cannot cycle and every row reaches row $0$, the one parentless row, with
an independent pointer-doubling pass confirming it; and that the rows are
in the canonical order, so that no reordering can pass unnoticed. Finally
the group element that its own canonicalization applied is read off as
$(\sigma_c,\varepsilon_c,s_c)$ and handed to a second, independent
implementation of the action --- one built from the defining formula, with
the inverse permutation and a reorientation-parity table rather than the
placement machinery --- which confirms the displayed identity exactly, as
$126$-bit sign vectors, for every one of the $9{,}276{,}594$ edges.

Validity of the mutant makes $B_{j_c}$ a mutable basis of the parent, and
the child is by construction a $G'$-translate of that mutant, so the two
classes really are adjacent in $\Gu^{9,4}$; the identity is the witness of
which translate it is, and a second implementation of the action confirms
it. With the tree structure, every listed class is therefore joined to the
root by a path of certified mutation edges.

This was run in full, not sampled: all $9{,}276{,}595$ classes
reconstructed and all $9{,}276{,}594$ edges verified, in $837$ seconds on
two processes, all $50$ checks passing --- $819$ of those seconds in the
reconstruction, the remaining $15$ in the integrity hashes, the structural
checks, the root, the sort that establishes distinctness and the mass
arithmetic. A reader without the legacy arrays runs the same thing minus
the five cross-checks of the third qualification below. Deriving the keys
rather than reading them costs nothing measurable: the $145$~MB artifact,
where they were given, took $853$ seconds on the same machine, since it
canonicalized each class only to compare against a stored key and then made
a second pass for the tree. The work is sharded and checkpointed, and the
rows of one depth are independent, which is what keeps it parallel although
each class is reconstructed from its parent.

On a $20{,}000$-row sub-certificate --- a prefix of the canonical order,
which is closed under the parent map and so is a complete certificate in
its own right --- twelve sabotages are each rejected: a duplicated class
(an edge re-pointed onto a class already listed), a corrupted root key, a
corrupted stabilizer histogram, a Grassmann--Pl\"ucker-invalid mutant, a
truncated certificate, a stale hash, a parent that does not precede its
row, a permutation of the parents preserving their multiset, a parent
re-pointed to another class at the same depth, a malformed gap stream, an
out-of-range mutated-basis index, and rows out of the canonical order.
Each is built at run time by tampering with a prefix of the real
certificate. Five carry a regenerated, internally consistent manifest ---
fresh hashes, and the count, mass and stabilizer histogram recomputed from
the sabotaged data --- so that only a substantive mathematical check can
catch them, and the three that re-point an edge are additionally asserted
at construction time to be semantically real: the tampered edge is
required to become invalid or to land on a class already listed, so that a
perturbation absorbed by a symmetry cannot produce a vacuous canary. A
sixth corrupts nothing but the manifest's stabilizer histogram. The
truncation deliberately leaves the totals unrepaired and the stale hash
leaves the manifest untouched, so that the count-and-mass arithmetic and
the integrity check are exercised too. The last four break the encoding,
the parent order or the row order; three of them are refused outright,
before any reconstruction starts, and the fourth --- two rows swapped out
of the canonical order --- is caught there and then caught again by the
mathematics downstream.

Three sabotages that the $145$~MB form admitted have no analogue here, and
the reason is the point of the new format: a corrupted stabilizer order, a
corrupted voltage and a second parentless class are not expressible,
because the certificate carries neither stabilizer orders nor voltages and
its root is structurally unique. Their nearest analogues --- a corrupted
histogram in the manifest, a corrupted mutated basis, and a parent that
does not precede its row --- are in the list above.

The certificate is tracked in git, so a reader who clones the repository
can check it where it lies, with no download and no regeneration; the
manifest pins it by the \texttt{SHA-256} of each of its arrays and of the
file itself, which is stored uncompressed with a fixed timestamp and so
reproduces byte for byte. The two legacy arrays are not tracked and are
not needed for any of the above. Regenerating them --- about four hours
for the search on four cores, six minutes for the export --- reproduces
the certificate exactly, and is what the cross-check of the third
qualification below requires.

\emph{Three qualifications, and they are not cosmetic.} First, ``extremal
in its own orbit'' is relative to the convention recorded in the manifest,
which maximises over the relabellings respecting an invariant colouring of
the ground set, not over all of $S_n$: the unrestricted maximum differs
from the reconstructed key on most classes and would cost some $160$
CPU-days here. It is still a well-defined function of the $G'$-orbit, the
colouring being invariant under reorientation and global negation and
equivariant under relabelling; and that --- not extremality as such --- is
what makes distinct keys certify distinct classes, and what makes the key
of a class computable from any representative of its orbit, which is what
the reconstruction uses. A reader who wants that fact rather than the word
``canonical'' must read the forty lines defining the
colouring. Second, the checker certifies the
left-hand side only. The target $1{,}722{,}704{,}635{,}330{,}560$ enters it
as a constant; an option re-adds the tracked $2{,}628$-row extension table
and confirms that it sums to that number, but nothing recomputes the
extension counts, nor re-derives the $(8,4)$ catalog they are taken over.
Those are cross-checked against every published count that exists, and the
identity closes exactly in the smaller cases where the answer is
independently known, but they are reproducible rather than certified.

Third, and this one is new with the compact certificate: the class keys
and the stabilizer orders are now \emph{derived} by the checker, not
compared against values the search recorded. That cuts both ways. A
certificate cannot misreport a key or a stabilizer order that it does not
contain, so a whole family of defects --- and of sabotages --- is no
longer expressible; but the earlier arrangement had two unrelated programs
agreeing, the generator's pruned canonicalization search and the checker's
exhaustive enumeration, and a reader who runs only what the repository
holds no longer sees that agreement. It can still be seen: the checker's
\texttt{--legacy-crosscheck} compares the reconstruction against the
arrays the search wrote, row by row, and on the untracked
$(9,4)$ arrays every one of the $9{,}276{,}595$ rows agrees in both its key
and its stabilizer order. That is corroboration of implementation agreement,
not a link in the chain: nothing in the argument above appeals to it.

The consequence deserves stating exactly, because the two halves of the
coverage claim now stand on different footings. That the $9{,}276{,}595$
listed classes are valid, pairwise inequivalent, carry their true
stabilizer orders and lie in a \emph{single} component of $\Gu^{9,4}$ is
certificate-backed outright, under the stated canonical convention, and
does not depend on the target at all: it is what the validity,
distinctness, stabilizer and spanning-tree checks above establish. That
they are \emph{all} the classes
is what the mass identity adds, and it is reproducible-only, since the
target is. So Theorem~\ref{thm:count}, and the last step of the coverage
half of Theorem~\ref{thm:main} --- from ``these $9{,}276{,}595$ classes
form one component'' to ``$\Gu^{9,4}$ is connected'' --- rest on a number
this project computes but does not certify. Nothing else in the coverage
half does.

\emph{One earlier observation was withdrawn.} An intermediate report from
this project described a structural phenomenon at $(9,4)$: the sign part
of the holonomy appeared to stay trivial over a large neighbourhood of
the alternating oriented matroid, which was read as evidence that the
labelled question there was delicate and perhaps negative. It was an
artifact of a defect in the search engine, which computed the stabilizer
of each newly discovered class and retained only its order, thereby
harvesting family (ii) of Lemma~\ref{lem:gens} but never family (i). With
the missing generators supplied the sign part is full immediately. The
observation is retracted. The defect escaped notice because the smaller
cases had run almost entirely through a different code path; the
regression suite now forces the path that the large runs take. We record
this because the failure mode is general: an incomplete harvest produces
a plausible \emph{negative} structural finding, and a negative finding is
exactly what one is predisposed to believe when a counterexample has been
conjectured.

\section*{Relation to \cite{KM23}}

Their rank-$3$ theorem $\Gb^{n,3}$ is theirs and is used here; our
labelled results add $(8,4)$, $(9,4)$ and $(9,5)$, which is what closes
$n\le 9$. Our $\Gm$ results for $n\le 9$ reprove theirs; the difference
is not a stronger theorem but that the computation here ships
certificates and checkers, whereas their connectivity computation ---
published in \cite{KM23}, with its method described there --- is reported
without publicly archived code, data or certificates. Their suspicion of
a labelled counterexample is not refuted in general; it is refuted in the
range $n\le 9$, which is where they had verified the other two levels.
Our Theorem~\ref{thm:count} confirms the value in their Table 1 against
\cite{FMM13}; the $(9,3)$ entry of that table needs the correction noted
in \S1.

\section*{Provenance}

This work was carried out by AI systems (Claude, with adversarial review
by Gemini and GPT models) directing exact computation under human
direction, with deliberately-failing controls in every search and
independent re-verification of the serialized certificates (with the
scope limits of \S\ref{sec:honesty}). No claim here rests on unverified
model output: the verdicts rest on those certificates and on the short
arguments above, both checkable by the reader. The author
directed the work and is responsible for the claims. The retraction in
\S\ref{sec:honesty} was found by this process, not despite it.

\begin{thebibliography}{9}
\bibitem{Finschi} L.~Finschi, \emph{Catalog of oriented matroids},
  \url{https://finschi.com/math/om/} (uniform rank-4 classes on 9
  elements computed by S.~Moriyama, December 2010).
\bibitem{FMM13} K.~Fukuda, H.~Miyata, S.~Moriyama,
  \emph{Complete enumeration of small realizable oriented matroids},
  Discrete Comput. Geom. \textbf{49} (2013), 359--381;
  \href{https://arxiv.org/abs/1204.0645}{arXiv:1204.0645}.
\bibitem{KM23} K.~Knauer, T.~Marc, \emph{Corners and simpliciality in
  oriented matroids and partial cubes}, European J. Combin. \textbf{112}
  (2023), 103714; \href{https://arxiv.org/abs/2002.11403}{arXiv:2002.11403}.
\bibitem{RS88} J.-P.~Roudneff, B.~Sturmfels, \emph{Simplicial cells in
  arrangements and mutations of oriented matroids},
  Geom. Dedicata \textbf{27} (1988), 153--170.
\end{thebibliography}

\end{document}
